\chapter{State of the Art}
\label{ch:related}

This chapter surveys the systems and protocols that address parts of the problem
\zkw tackles, and positions the present work relative to them. It is organised by
sub-problem: first the established whistleblowing platforms, then dead-man's-switch
implementations, decentralised key-management networks, and provenance
technologies. It closes with a critical look at the project's own earlier
TEE-based design and a gap analysis that motivates the architecture of
Chapter~\ref{ch:architecture}.

\section{Whistleblowing and Secure Submission Platforms}
\label{sec:rel-platforms}

The most widely deployed tools for receiving sensitive disclosures are
\emph{SecureDrop}~\cite{securedrop} and \emph{GlobaLeaks}~\cite{globaleaks}.
SecureDrop, maintained by the Freedom of the Press Foundation, is an open-source
submission system run by individual news organisations: a source connects over
Tor to an organisation's instance and uploads documents, which journalists
retrieve from an air-gapped viewing station. GlobaLeaks is a more general
software framework for building anonymous reporting initiatives, used by
anti-corruption NGOs and public bodies.

These platforms are mature and have a strong operational-security track record,
and they solve a genuinely hard problem: getting material from a source to a
journalist without exposing the source's network identity. They are, however,
\emph{custodial} and \emph{organisation-bound} by design. Each instance is run by
a specific organisation that operates the servers, holds the submitted material,
and is the locus of trust. This produces three limitations that \zkw is built to
overcome. First, the organisation is a single point of failure: it can be
compromised, subpoenaed, or shut down, and its operational security is only as
good as its staff and infrastructure. Second, these systems address the
\emph{transmission} of evidence, not the \emph{vulnerable interval}---they offer
no mechanism to release evidence automatically if a source is silenced before
they choose to submit. Third, they offer no notion of \emph{verifiable
provenance} for an anonymous source and no mechanism for \emph{anonymous reward};
credibility is established, if at all, through out-of-band journalistic
verification.

WikiLeaks popularised the use of an ``insurance file''---a large encrypted blob
released publicly in advance, with the decryption key withheld and released only
if something happened to the organisation. This is a manual, all-or-nothing dead
man's switch whose key custody rests entirely on a single human decision. It
illustrates both the demand for the mechanism this thesis automates and the
fragility of doing it by hand.

\section{Dead Man's Switch Implementations}
\label{sec:rel-dms}

Digital dead-man's-switch services have existed for years in centralised form:
``email this for me if I don't log in for 30 days'' products that depend entirely
on the operator's honesty and survival. They reproduce the custodial problem in
its purest form and are out of scope as serious protections.

On-chain implementations are more interesting. \emph{Sarcophagus}~\cite{sarcophagus}
is a decentralised dead-man's-switch protocol on Ethereum that uses Arweave for
permanent storage---a design this thesis draws inspiration from. In Sarcophagus, a
user (the ``embalmer'') encrypts a file in two layers: an inner layer for the
ultimate recipient and an outer layer for a staked node operator (an
``archaeologist''). If the user fails to ``re-wrap'' (check in) before a
resurrection time, the archaeologist removes the outer layer, publishing the
inner-encrypted file for the recipient to decrypt. Sarcophagus pioneered the
combination of on-chain trigger logic with permanent storage, but it carries two
drawbacks for the present use case. Its security relies on an \emph{economic}
incentive model with its own token (\$SARCO): the user must keep paying node
operators, and a lapse in payment can cause the switch to fail or the data to be
lost. And its proxy-re-encryption-style flow is oriented toward a \emph{known
recipient} rather than a public or as-yet-unknown future journalist.

Hackathon prototypes such as ETHGlobal's ``Deadman Switch''~\cite{ethglobalDeadman}
demonstrate the heartbeat pattern---an on-chain timestamp updated by a periodic
\texttt{checkIn()} and an \texttt{isDeceased()} predicate---which \zkw adopts and
generalises. What these prototypes typically lack is a credible answer to key
custody (they often hand-wave the decryption step) and any treatment of
provenance or anonymous reward.

\zkw differs from all of the above in that it decouples the four concerns---trigger,
key custody, provenance, and reward---and assigns each to a purpose-built
decentralised primitive, rather than bundling them into a single token-incentivised
protocol or leaving key custody unspecified.

\section{Decentralised Key-Management Networks}
\label{sec:rel-keymgmt}

The hardest sub-problem is conditional key custody: who holds the decryption key,
and what convinces them to release it. Three families of solution exist.

\textbf{Trusted execution environments}, discussed in Section~\ref{sec:bg-tee},
seal the key in hardware. This is the approach of the project's V1 and is rejected
for the reasons given there and revisited in Section~\ref{sec:rel-v1}.

\textbf{Proxy re-encryption (PRE) networks}, exemplified by the \emph{Threshold
Network} (the merger of NuCypher and Keep)~\cite{thresholdNetwork, nucypher},
distribute re-encryption keys across nodes via secret sharing. The data owner
encrypts to their own public key; when access is granted, the network transforms
the ciphertext so that it is decryptable by a designated recipient, without any
node ever seeing the plaintext. PRE is mathematically elegant and genuinely
non-custodial, but it is fundamentally \emph{point-to-point}: the recipient must
generally be known, with a public key, \emph{at encryption time}. A dead man's
switch for whistleblowing frequently needs to release to the \emph{public} or to
an unknown future journalist, which fits PRE poorly.

\textbf{Programmable threshold key management}, exemplified by the \emph{Lit
Protocol}~\cite{litprotocol}, is the family this thesis adopts. Lit is a
decentralised network whose nodes hold key shares (using threshold BLS) and that
exposes two relevant capabilities. \emph{Access Control Conditions} (ACCs) let
data be encrypted such that the network will only release the decryption shares
when a specified condition holds---including arbitrary on-chain state, such as
the return value of a contract call. \emph{Lit Actions} are immutable JavaScript
programs that can express richer release logic. For a dead man's switch, the ACC
simply requires that \texttt{DeadMansSwitch.isDeceased(user)} return true. Unlike
PRE, this naturally supports release to the public (the condition, not a specific
recipient's key, gates access); unlike a TEE, the key is split across a network
and its security is threshold-cryptographic. Lit's developer-facing SDK abstracts
the MPC so that an application can encrypt and decrypt against conditions without
implementing the threshold protocol itself, which is what makes it usable as the
key-management layer of this project.

It must be noted---and Chapter~\ref{ch:evaluation} treats this directly---that
Lit's own technology has evolved during the lifetime of this work. The network
generation used here (Naga, via the v8 SDK) preserves the multi-party-computation
model the thesis argues for, whereas a newer Lit runtime reintroduces TEE-based
execution. The design deliberately targets the MPC path.

\section{Provenance and Reputation Technologies}
\label{sec:rel-provenance}

Establishing the credibility of an anonymous claim without identity is a newer
capability. As described in Section~\ref{sec:bg-zk}, \emph{Reclaim
Protocol}~\cite{reclaim} uses zkTLS to turn any authenticated HTTPS session into a
selective, zero-knowledge proof of its contents. This is the first widely
available technology that lets a source prove ``I am an active employee of
organisation X'' or ``I control this account'' without revealing their identity,
and it is the provenance layer adopted here. \emph{ZK-Email}~\cite{zkemail}
provides a complementary capability by proving properties of DKIM-signed emails in
zero knowledge---for instance, that one received an email from a specific
corporate domain---which is a natural future extension.

Earlier reputation systems on-chain were either identity-bound (defeating
anonymity) or based on token holdings and social graphs (not evidence of
insider access). Reclaim and ZK-Email are distinctive in producing
\emph{verifiable, identity-free} evidence of a real-world relationship, which is
exactly the signal a journalist needs to triage anonymous tips.

\section{The V1 Design and Its Limitations}
\label{sec:rel-v1}

The immediate predecessor of this work used Intel SGX, accessed through the iExec
decentralised-computing network, to hold decryption keys for the dead man's
switch. The intent was sound---move key custody off any single human or
server---but the realisation embeds the trust in hardware. As
Section~\ref{sec:bg-tee} detailed, this is vulnerable to side-channel and
transient-execution attacks (Foreshadow, SGAxe), depends on a single chip
vendor's attestation infrastructure as a centralised root of trust, and ties
liveness to a narrow set of specialised nodes. Comparative analyses of
confidential-computing technologies reach the same
conclusion~\cite{phalaTeeMpc}: TEEs offer excellent performance and a simple
programming model, but for adversaries capable of physical or supply-chain
attacks they are a weaker guarantee than distributed cryptography. The decision to
re-architect around threshold cryptography is the direct response to this
analysis and the central thesis of this work.

\section{Gap Analysis and Positioning}
\label{sec:rel-gap}

Table~\ref{tab:rel-gap} summarises the landscape against the four capabilities
that a complete confidential-whistleblowing system requires: an autonomous
release trigger, non-hardware key custody, identity-free provenance, and
unlinkable reward.

\begin{table}[htbp]
\centering
\caption{Capability comparison of related systems.}
\label{tab:rel-gap}
\small
\begin{tabular}{@{}lcccc@{}}
\toprule
\textbf{System} & \textbf{Auto.} & \textbf{Non-HW key} & \textbf{ZK} & \textbf{Unlinkable} \\
 & \textbf{trigger} & \textbf{custody} & \textbf{provenance} & \textbf{reward} \\
\midrule
SecureDrop / GlobaLeaks      & No  & N/A (custodial)      & No  & No  \\
WikiLeaks insurance file     & Manual & No (human)        & No  & No  \\
Sarcophagus                  & Yes & Economic / staked    & No  & No  \\
ETHGlobal DMS prototypes     & Yes & Unspecified          & No  & No  \\
Threshold Network (PRE)      & Partial & Yes (point-to-point) & No & No \\
Project V1 (iExec / SGX)     & Yes & No (hardware TEE)    & No  & No  \\
\textbf{\zkwns (this work)}  & \textbf{Yes} & \textbf{Yes (MPC/TSS)} & \textbf{Yes} & \textbf{Yes} \\
\bottomrule
\end{tabular}
\end{table}

No existing system addresses all four capabilities together. SecureDrop and
GlobaLeaks excel at secure transmission but are custodial and address neither the
vulnerable interval nor provenance nor reward. Sarcophagus and the hackathon
prototypes provide an autonomous trigger but with economic or unspecified key
custody and no provenance or reward. PRE networks provide non-custodial keys but
are point-to-point and silent on the other three concerns. The V1 design provides
an autonomous trigger and decentralised \emph{computation} but anchors trust in
hardware.

The contribution of \zkw is to occupy the final row of
Table~\ref{tab:rel-gap}: to combine an autonomous on-chain trigger, threshold
(non-hardware) key custody, zero-knowledge provenance, and unlinkable reward into
a single, internally consistent architecture, with a careful specification of the
trust boundaries between the independent networks that provide each capability.
The next chapter sets out that architecture.
